% MLR-X User Manual
\documentclass[11pt,a4paper]{article}
\usepackage[utf8]{inputenc}
\usepackage[T1]{fontenc}
\usepackage[english]{babel}
\usepackage{hyperref}
\usepackage{graphicx}
\usepackage{longtable}
\usepackage{array}
\usepackage{enumitem}
\usepackage{geometry}
\geometry{margin=2.5cm}

\title{MLR-X 1.0\\User Guide and Technical Reference}
\author{MLR-X Project}
\date{November 2025}

\begin{document}
\maketitle
\tableofcontents
\newpage

\section{Overview}
MLR-X 1.0 is a multiple linear regression toolkit with both graphical (GUI) and command-line (CLI) modes. It is designed for small and large datasets alike, offering optimized routines, parallel execution, and interactive diagnostic plots.\footnote{The public header announces the product version and tagline.}

\begin{figure}[h!]
  \centering
  \includegraphics[width=0.9\textwidth]{fig_portada_placeholder}
  \caption{MLR-X start screen (reference illustration).}
  \label{fig:portada}
\end{figure}

\subsection{Licensing and authorship}
\begin{itemize}[leftmargin=1.4cm]
  \item License: GNU Affero General Public License v3.0 (AGPL-3.0).
  \item Copyright: 2025, Jackson J. Alcázar.
  \item Official repository: \url{https://github.com/Jacksonalcazar/MLR-X}.
  \item Donations: \url{https://www.paypal.com/donate?business=jacksonalcazar@gmail.com\&currency_code=USD}.
\end{itemize}

\subsection{Versioning}
This manual documents version \textbf{1.0}. The executable accepts the \verb|--version| flag to print the version from the command line. The file \texttt{version\_info.txt} defines both file and product version numbers as \texttt{1.0.0.0} to keep shipped binaries consistent.\footnote{The \texttt{VERSION} constant in the source code is set to 1.0.}

\section{Key features}
\begin{itemize}[leftmargin=1.2cm]
  \item \textbf{Scalability and efficiency}: handles small and large datasets through deferred loading of heavy libraries and automatic parallel execution.\cite{mlrx-index}
  \item \textbf{GUI and CLI modes}: an intuitive GUI for quick exploration and a console utility for reproducible pipelines.\cite{mlrx-index}
  \item \textbf{Reports and diagnostics}: standard metrics (MAE, MSE, R\textsuperscript{2}) and export-ready reports.\cite{mlrx-index}
  \item \textbf{Open source}: distributed under AGPL-3.0 with a public repository for collaboration.\cite{mlrx-index}
  \item \textbf{Official downloads}: packages for Linux, Windows (GUI and GUI+CLI), and an announced macOS build.\cite{mlrx-index}
\end{itemize}

\section{Installation}
\subsection{Packaged downloads}
The official site provides ready-to-use binaries:
\begin{itemize}
  \item Linux (direct download link).
  \item Windows GUI.
  \item Windows GUI + CLI.
  \item macOS (marked as ``coming soon'').
\end{itemize}
Replace the example links with the final URLs from the downloads page.\cite{mlrx-index}

\subsection{System requirements}
\begin{itemize}
  \item \textbf{Operating system}: Windows, Linux, or macOS.
  \item \textbf{Python}: the source uses Python 3 with dependencies such as NumPy, pandas, scikit-learn, statsmodels, and Matplotlib, loaded on demand to optimize startup.\cite{mlrx-source-imports}
  \item \textbf{Graphical interface}: the GUI relies on Tkinter; ensure the corresponding system libraries are available.
  \item \textbf{Resources}: the application will use multiple CPU cores when available.
\end{itemize}

\section{Using the graphical application (GUI)}
\subsection{Launch}
Running the program without arguments opens the main GUI. Heavy scientific modules are preloaded in the background to reduce latency during the first operations.\cite{mlrx-source-main}

\subsection{Basic workflow}
\begin{enumerate}
  \item \textbf{Load data}: select a file (semicolon-delimited by default). You can specify the dependent variable (by name or as the last column) and mark non-variable columns.
  \item \textbf{Configure the model}: choose the search method (EPRS-S or traditional ``all subsets''), the maximum number of predictors, the number of seeds, and seed size.
  \item \textbf{Validation}: set the covariance matrix type (HC0--HC3 or non-robust) and thresholds for collinearity (correlation and VIF).
  \item \textbf{Run}: start the fitting process; the interface shows progress and enables inspection of primary metrics.
  \item \textbf{Visualization and diagnostics}: access residual plots, leverage plots, Cook's distances, and other diagnostics; axes are customizable and points can be labeled by split or ID.\footnote{Insert GUI captures at markers such as \texttt{fig:residuos}, \texttt{fig:cooks}, etc.}
  \item \textbf{Export}: save plots to PNG, TIFF, PDF, or SVG, and export result tables to CSV from the GUI view.
\end{enumerate}

\subsection{Accessibility and usability}
\begin{itemize}
  \item Controls to toggle legends and gridlines on plots.
  \item Axis labels and font sizes can be edited directly from the interface.
  \item Marker shapes and colors can be customized per dataset when comparing splits.
\end{itemize}

\section{Command-line usage (CLI)}
\subsection{Basic invocation}
\begin{verbatim}
$ python MLRX_source.py --version
MLR-X version 1.0

$ python MLRX_source.py path/to/config.conf
\end{verbatim}
The first command prints the version. The second loads a configuration file and runs the full console workflow.\cite{mlrx-source-main}

\subsection{Configuration file}
The CLI reads a text file structured into blocks and key-value pairs. Each section can start with a header (for example, \# Dataset) followed by options. Fields are validated and normalized during loading.\cite{mlrx-source-config}

\paragraph{Core fields}
\begin{longtable}{>{\bfseries}p{4cm}p{10cm}}
Data path & Path to the input CSV.\\
Delimiter & Column separator (\verb|;| by default, \verb|,|, tab, or a literal text value).\\
Dependent choice & Target variable (\verb|last| or a column name).\\
Non variable spec & Columns to exclude from the predictor space.\\
Exclude constant & Drops constant columns based on a variance threshold.\\
Max vars & Maximum number of predictors per model.\\
N seeds / Seed size & Parameters for the EPRS-S search.\\
Cov type & Covariance matrix type (\texttt{nonrobust}, \texttt{HC0}, \texttt{HC1}, \texttt{HC2}, \texttt{HC3}).\\
Signif lvl & Significance level for filtering.\\
Corr threshold / VIF threshold & Limits for detecting collinearity.\\
Report R2 threshold & Cutoff to include models in reports.\\
Export limit & Maximum number of models in the CSV export.\\
N jobs & Cores to use (normalized to at least 1).\\
Iterations mode & \texttt{auto}, \texttt{manual} (requires \textit{Max iterations per seed}) or \texttt{converge}.\\
Clip predictions & Enables clipping intervals (\textit{clip\_low}, \textit{clip\_high}).\\
Split mode & \texttt{none}, \texttt{random} (test percentage), \texttt{manual} (train/test IDs) or \texttt{external} (path and optional delimiter).\\
Output path & Path to the results CSV.\\
Target metric & Objective metric (R\textsuperscript{2} or variants).\\
Method & \texttt{eprs} or \texttt{all\_subsets}.\\
\end{longtable}

\paragraph{Selecting options}
Blocks with lists (for example, delimiters or methods) can be declared via numeric indices using a line such as ``\texttt{selected option: <n>}'' inside the corresponding block. The system validates that indices are in range and that options exist.\cite{mlrx-source-config}

\subsection{Execution and progress}
After loading the configuration, the CLI reports basic dataset info, available predictors, the chosen method, and key parameters. Progress is reported by stage (search or VIF calculation) and may print inline updates while running.\cite{mlrx-source-cli}

\subsection{Output and export}
Evaluated models are exported to a CSV with normalized columns (number of variables, R\textsuperscript{2}, RMSE, MAE, and validation variants such as LOO, k-fold, and external). Variable names are normalized and process metadata is included (CPU time, number of models found and explored, split configurations).\cite{mlrx-source-export}

\section{Diagnostics and plots}
The GUI offers visualizations such as:
\begin{itemize}
  \item Residuals vs. leverage plot with reference lines and a configurable leverage threshold.
  \item Cook's distance (including LOO variant) with \(4/n\) references and guide lines.
  \item Histograms and QQ-plots to assess residual normality (configured in the source code).
\end{itemize}
Each figure should be cited in the text when inserted in Overleaf (for example, Figure~\ref{fig:cooks}).

\section{Best practices}
\begin{itemize}
  \item Validate CSV delimiters and headers before loading.
  \item Adjust \texttt{n\_jobs} based on CPU availability to balance time and resources.
  \item Use \texttt{iterations mode = converge} to let the EPRS-S algorithm set iterations by convergence.
  \item Review collinearity using the correlation and VIF thresholds before accepting a model.
  \item Export results with a reasonable \texttt{export\_limit} to keep files manageable.
\end{itemize}

\section{References and citation}
For academic work, cite the software as:\newline
\textit{Alcázar, Jackson J. ``MLR-X 1.0: A Scalable Software for Multiple Linear Regression on Small and Large Datasets''.} A BibTeX entry is included in the distribution with the key \texttt{alcazar\_mlr\_x\_1\_0}.\cite{mlrx-source-citation}

\section{Support and contributions}
\begin{itemize}
  \item \textbf{Repository}: \url{https://github.com/Jacksonalcazar/MLR-X}
  \item \textbf{Issues or improvements}: open an issue in the official repository.
  \item \textbf{PDF manual}: available in the repository (replace the example link if the version changes).
  \item \textbf{Contact}: author's email (see footer on the official site).
\end{itemize}

\section{Appendix: Minimal configuration example}
\begin{verbatim}
# Dataset
data_path: data.csv
delimiter: ;
dependent_choice: last
non_variable_spec: ID
exclude_constant: true
constant_threshold: 90

# Method
method: eprs
max_vars: 10
n_seeds: 20
seed_size: 4
cov_type: HC3
signif_lvl: 0.05
corr_threshold: 0.90
vif_threshold: 4.0
report_r2_threshold: 0.80
iterations_mode: auto
export_limit: 500
n_jobs: 4

# Validation
split_mode: random
random_test_size_percent: 20

# Output
output_path: results/models.csv

target_metric: R2
clip_enabled: false
\end{verbatim}

\section{Version history}
\begin{itemize}
  \item \textbf{1.0 (2025)}: First public release with GUI and CLI, EPRS-S search, CSV export, and the initial PDF manual.
\end{itemize}

\bibliographystyle{plain}
\begin{thebibliography}{9}
\bibitem{mlrx-index} Public MLR-X site section describing features and downloads.\\
\url{https://example.com/replace-with-final-url}

\bibitem{mlrx-source-imports} Source code of MLRX\_source.py defining deferred imports and version.\\
\url{https://example.com/replace-with-final-url}

\bibitem{mlrx-source-config} Configuration parsing code in MLRX\_source.py.\\
\url{https://example.com/replace-with-final-url}

\bibitem{mlrx-source-cli} \texttt{run\_cli} entry point and progress handling in MLRX\_source.py.\\
\url{https://example.com/replace-with-final-url}

\bibitem{mlrx-source-export} CSV export logic in MLRX\_source.py.\\
\url{https://example.com/replace-with-final-url}

\bibitem{mlrx-source-citation} Citation text and BibTeX entry in MLRX\_source.py.\\
\url{https://example.com/replace-with-final-url}
\end{thebibliography}

\end{document}
